% Part: first-order-logic
% Chapter: syntax-and-semantics
% Section: free-vars-sentences

\documentclass[../../../include/open-logic-section]{subfiles}

\begin{document}

\olfileid[pt]{fol}{syn}{fvs}

\olsection{\printtoken{P}{variable} Livres e \printtoken{P}{sentence}}

\begin{defn}[Ocorrências livres de !!a{variable}]
\ollabel{defn:free-occ}
As ocorrências \emph{livres} de !!a{variable} em !!a{formula} são definidas
indutivamente como segue:
\begin{enumerate}
\item \indcase*{!A}{$!A$ é atômica}{todas as ocorrências de !!{variable} em
  $\indfrm$ são livres.}

\tagitem{prvNot}{\indcase{!A}{\lnot !B}{as ocorrências livres de !!{variable}
  de $\indfrm$ são exatamente as de $!B$.}}{}

\item \indcase{!A}{(!B \ast !C)}{as ocorrências livres de
  !!{variable} de $\indfrm$ são as em $!B$
  juntamente com as em~$!C$.}

\tagitem{prvAll}{\indcase{!A}{\lforall[x][!B]}{as ocorrências livres de !!{variable}
  em $\indfrm$ são todas as em~$!B$ exceto as
  ocorrências de~$x$.}}{}

\tagitem{prvEx}{\indcase{!A}{\lexists[x][!B]}{as ocorrências livres de !!{variable}
  em $\indfrm$ são todas as em~$!B$ exceto as
  ocorrências de~$x$.}}{}
\end{enumerate}
\end{defn}

\begin{defn}[Variáveis Ligadas]
Uma ocorrência de !!a{variable} em uma fórmula~$!A$ é \emph{ligada} se
não é livre.
\end{defn}

\begin{prob}
Dê uma definição indutiva das ocorrências de variáveis ligadas nos
moldes de \olref[fol][syn][fvs]{defn:free-occ}.
\end{prob}

\begin{defn}[Escopo]
\iftag{prvAll}{Se $\lforall[x][!B]$ é uma ocorrência de uma subfórmula
  em uma fórmula~$!A$, então a ocorrência correspondente de~$!B$ em~$!A$
  é chamada de \emph{escopo} da ocorrência correspondente
  de~$\lforall[x]$. \iftag{prvEx}{Analogamente para $\lexists[x]$.}{}}{Se
  $\lexists[x][!B]$ é uma ocorrência de uma subfórmula em uma
  fórmula~$!A$, então a ocorrência correspondente de~$!B$ em~$!A$ é
  chamada de \emph{escopo} da ocorrência correspondente
  de~$\lexists[x]$.}

Se $!B$ é o escopo de uma ocorrência de quantificador
\iftag{prvAll}{$\lforall[x]$\iftag{prvEx}{ ou
    $\lexists[x]$}{}}{$\lexists[x]$} em~$!A$, então as ocorrências livres de
$x$ em~$!B$ são ligadas em \iftag{prvAll}{$\lforall[x][!B]$\iftag{prvEx}{ e
$\lexists[x][!B]$}{}}{$\lexists[x][!B]$}. Dizemos que essas
ocorrências são \emph{ligadas pela}
ocorrência de quantificador mencionada.
\end{defn}

\begin{ex}
Considere a seguinte fórmula:
\[
\lexists[\Obj v_0][\underbrace{\Atom{\Obj A^2_0}{\Obj v_0,\Obj v_1}}_{!B}] 
\]
$!B$ representa o escopo de $\lexists[\Obj v_0]$. 
O quantificador liga a ocorrência de $\Obj v_0$ em $!B$, mas
não liga a ocorrência de $\Obj v_1$. Então $\Obj v_1$ é
uma variável livre neste caso.


Agora podemos ver como isso poderia funcionar em uma 
!!{formula}~$!A$ mais complicada:
\[
\lforall[\Obj v_0][\underbrace{(\Atom{\Obj A^1_0}{\Obj v_0} \lif
    \Atom{\Obj A^2_0}{\Obj v_0, \Obj v_1})}_{!B}] \lif \lexists[\Obj
  v_1][\underbrace{(\Atom{\Obj A^2_1}{\Obj v_0, \Obj v_1} \lor \lforall[\Obj v_0][\overbrace{\lnot \Atom{\Obj A^1_1}{\Obj v_0}}^{!D}])}_{!C}]
\]
$!B$ é o escopo do primeiro $\lforall[\Obj v_0]$, $!C$ é o escopo
de $\lexists[\Obj v_1]$, e $!D$ é o escopo do segundo
$\lforall[\Obj v_0]$.  O primeiro $\lforall[\Obj v_0]$ liga as
ocorrências de $\Obj v_0$ em~$!B$, $\lexists[\Obj v_1]$ liga a ocorrência
de $\Obj v_1$ em $!C$, e o segundo $\lforall[\Obj v_0]$ liga a
ocorrência de $\Obj v_0$ em~$!D$.  A primeira ocorrência de $\Obj v_1$
e a quarta ocorrência de $\Obj v_0$ são livres em~$!A$. A última
ocorrência de $\Obj v_0$ é livre em $!D$, mas ligada em $!C$ e~$!A$.
\end{ex}

\begin{defn}[Sentença]
!!^a{formula}~$!A$ é \article{sentence} \emph{!!{sentence}} sse
não contém ocorrências livres de !!{variable}s.
\end{defn}

% add examples!

\end{document}
